\documentclass[11pt,a4paper]{article}
\usepackage[margin=25mm]{geometry}
\usepackage{fontspec}
\defaultfontfeatures[\rmfamily,\sffamily]{Ligatures=Common}
\setmainfont{TeX Gyre Pagella}
\setmonofont{DejaVu Sans Mono}[Scale=0.80]
% The monospaced font lacks these two symbols used by Lean's notation.
% XeTeX changes font for the original Unicode characters, preserving copy/paste.
\newfontfamily\leanSymbolFont{FreeSerif}[Scale=0.80]
\newXeTeXintercharclass\LeanSymbols
\XeTeXcharclass"2964=\LeanSymbols
\XeTeXcharclass"1D4DD=\LeanSymbols
\XeTeXinterchartoks 0 \LeanSymbols={\begingroup\leanSymbolFont}
\XeTeXinterchartoks \LeanSymbols 0={\endgroup}
\XeTeXinterchartoks 4095 \LeanSymbols={\begingroup\leanSymbolFont}
\XeTeXinterchartoks \LeanSymbols 4095={\endgroup}
\XeTeXinterchartokenstate=1
\usepackage{amsmath,amssymb,amsthm}
\usepackage{booktabs}
\usepackage{xcolor}
\definecolor{linkblue}{HTML}{245878}
\usepackage{fvextra}
\fvset{fontsize=\footnotesize,breaklines=true,breakanywhere=true,
  breakautoindent=false,breakindent=0pt}
\usepackage{xurl}
\usepackage[colorlinks=true,linkcolor=linkblue,urlcolor=linkblue]{hyperref}
\setlength{\emergencystretch}{3em}
\newtheorem{claim}{Claim}[section]
\newcommand{\formal}[1]{\par\smallskip\noindent
  \textit{Extracted Lean statement: Appendix~\ref{decl:#1}.}\par\smallskip}
\newcommand{\scope}{\par\smallskip\noindent\textbf{Scope.} }
% Generated by proof_companion/render.sh; do not edit.
\newcommand{\CompanionTheoremCount}{12}
\newcommand{\CompanionModuleCount}{4}
\newcommand{\CompanionTacticRecordCount}{1357}

\title{Proof companion\\\large Coupling-gated cortical coherence}
\author{A companion to the repository's Lean development}
\date{First coverage stage: 16 September 2026}

\begin{document}
\maketitle
\begin{abstract}
This document presents mathematical explanations of
\CompanionTheoremCount{} selected Lean results from
\CompanionModuleCount{} source modules. It places the assumptions and the
argument of each result beside an appendix generated from Lean's elaborated
statements. The saved extraction also contains
\CompanionTacticRecordCount{} nested tactic records, their before/after goal
states, proof dependencies, and the fields of project structures reached through
statement types. The mathematical prose is an editorial draft checked against
the source during preparation; Lean verifies the formal proofs, not this prose.
\end{abstract}

\tableofcontents

\section{Reading and auditing the companion}
The first coverage stage follows four kinds of argument: exact and approximate
gluing, the critical exponent of stationary coherent solutions, the conditional
composition of the framework, and a concrete nonconstant family. The selection
is explicit in \nolinkurl{selection.json}. Each of its declarations has an
extracted statement in the appendix, including implicit type and instance
parameters. Definitions and hypotheses abbreviated in the mathematics are
spelled out where they carry substantive content.

The prose groups rewrites and elementary algebra into mathematical steps.
Full tactic histories remain in the JSON snapshots under \nolinkurl{data/};
Section~\ref{sec:witness-trace} displays a small actual trace. A nested tactic
record is an elaboration event, so its count includes wrappers and macro
expansions. It is not a count of independent mathematical steps.

Three distinctions guide the reading. A theorem can depend on physical
conditions supplied as structure fields. A formally valid conditional theorem
can have hypotheses that require separate witnesses. A dependency on a named
definition does not establish an empirical interpretation of that definition.
The examples below keep these distinctions visible.

The extractor rejects elaboration errors and checks that each selected result's
axioms are drawn only from \texttt{propext}, \texttt{Classical.choice}, and
\texttt{Quot.sound}. This is a check of logical dependencies. It does not remove
the hypotheses in the statement or the obligations carried by its structures.

\section{Exact and approximate gluing}\label{sec:gluing}

\subsection{The sheaf used by the development}
Let $X$ be a topological space with the measurable and Borel structure required
by the formal statement. Let $F_0$ assign to an open subset its finite spatial
measures, with restriction maps. The Lean definition
\texttt{probabilityPresheaf} denotes the sheafification $F$ of $F_0$.
Despite that identifier, the initial local data are unnormalised finite
measures; mass one is not part of this definition.

\begin{claim}
$F$ is a sheaf.
\end{claim}
\begin{proof}
Sheafification returns a sheaf together with its sheaf-property proof.
Projecting that proof establishes the claim for $F$.
\end{proof}
\formal{sheaf-property}
\scope This result uses the sheafification construction. It supplies no
physical mechanism for producing compatible local data.

\subsection{Compatible local sections have a unique extension}
\begin{claim}
Let $F$ be any sheaf of types on $X$, and let $(U_i)_{i\in I}$ be an open
cover of $X$. Suppose $s_i\in F(U_i)$ and
\[
 s_i|_{U_i\cap U_j}=s_j|_{U_i\cap U_j}\qquad(i,j\in I).
\]
Then exactly one $g\in F(X)$ satisfies $g|_{U_i}=s_i$ for every $i$.
\end{claim}
\begin{proof}
The unique-gluing form of the sheaf condition gives a section on
$\bigcup_i U_i$, with the specified restrictions, unique among sections there.
The cover equality identifies that union with $X$. Transport the section along
this equality. Functoriality of the restriction maps shows that its restriction
to each $U_i$ is $s_i$.

Conversely, transport any competing section on $X$ back to the union. Its
restrictions are again the $s_i$, so the uniqueness supplied by the sheaf
condition identifies it with the first section. Transporting forward gives the
required equality on $X$. In the Lean proof, equality of the relevant inclusion
maps uses the fact that morphisms between opens form a subsingleton.
\end{proof}
\formal{sheaf-gluing}

\begin{claim}
The same conclusion holds for the development's sheafified spatial measures.
\end{claim}
\begin{proof}
Apply the preceding claim to $F=\texttt{probabilityPresheaf}$ and use the
sheaf property established above. The compatible family and cover equality
are exactly the hypotheses of this application.
\end{proof}
\formal{probability-gluing}
\scope This proof is a direct term applying a previous theorem. It has no
tactic history; the extracted statement and proof dependencies still cover it.

\subsection{What equilibrium contributes}
A thermodynamic cover supplies local sections $s_i$, phases $\theta_i$, and a
finite, symmetric, strictly positive coupling matrix. Its structure includes
two substantive conditions: the phase configuration globally minimises the
reduced Kuramoto potential, and
\[
 \cos(\theta_i-\theta_j)=1
 \quad\Longrightarrow\quad
 s_i|_{U_i\cap U_j}=s_j|_{U_i\cap U_j}.
\]
The second condition is an obligation on the supplied local data.

\begin{claim}
A thermodynamic cover determines a unique global section with its prescribed
local restrictions.
\end{claim}
\begin{proof}
The theorem characterising global minima of the reduced Kuramoto potential
implies phase locking for the cover: $\cos(\theta_i-\theta_j)=1$ for every pair.
The overlap-agreement field then makes the local sections compatible. Apply
the spatial-measure gluing theorem to the cover and that compatible family.
\end{proof}
\formal{thermodynamic-gluing}
\scope This argument consumes the supplied equilibrium and overlap conditions.
The existence of a physical system satisfying them is a separate matter.

\subsection{Approximate agreement gives a controlled selection}
Let the site set and patch index set be finite. Let $s_i(x)\geq0$ be local
mass profiles represented as functions on the ambient sites. Let
$w_i(x)\geq0$ satisfy $w_i(x)=0$ outside $U_i$ and
$\sum_i w_i(x)=1$ at every site. Assume that for $x\in U_i\cap U_j$,
$|s_i(x)-s_j(x)|\leq\varepsilon$, where $\varepsilon\geq0$.

\begin{claim}
The selected profile $g(x)=\sum_j w_j(x)s_j(x)$ satisfies
$\sup_{x\in U_i}|g(x)-s_i(x)|\leq\varepsilon$ for every patch $i$.
\end{claim}
\begin{proof}
Fix $x\in U_i$. Every nonzero weight $w_j(x)$ has $x\in U_j$, so its profile
is within $\varepsilon$ of $s_i(x)$. Normalisation and the triangle inequality
therefore give
\[
 |g(x)-s_i(x)|
 =\left|\sum_j w_j(x)(s_j(x)-s_i(x))\right|
 \leq\sum_j w_j(x)\varepsilon=\varepsilon.
\]
The terms with zero weight contribute zero. The uniform metric on the finite
patch converts these pointwise bounds into the displayed profile-distance
bound. Lean isolates the weighted inequality in its helper
\texttt{weighted\_dist\_le} and then applies it at each site.
\end{proof}
\formal{approximate-gluing}
\scope The profiles need not have total mass one. The result provides an
approximation in the specified uniform mass metric, and assumes the overlap
bounds and the subordinate normalised weights.

\section{The stationary critical exponent}
Write
\[
 E(a)=\frac{\int_0^{2\pi}\sin^2\theta\,e^{a\cos\theta}\,d\theta}
             {\int_0^{2\pi}e^{a\cos\theta}\,d\theta}.
\]
This is \texttt{vonMisesSRatio}. Its denominator is strictly positive.

\subsection{The second-order expansion}
\begin{claim}
As $a\to0$,
\[
 E(a)=\frac12-\frac{a^2}{16}+o(a^2).
\]
\end{claim}
\begin{proof}
The previously proved differentiability results give $E\in C^2$,
$E(0)=1/2$, $E'(0)=0$, and $E''(0)=-1/8$. Taylor's theorem with a little-o
remainder yields
$E(a)=E(0)+E'(0)a+E''(0)a^2/2+o(a^2)$.
Substituting those values gives the claimed polynomial. The Lean proof applies
the general Taylor theorem, identifies its second iterated derivative using
the derivative lemmas, and simplifies its Taylor polynomial.
\end{proof}
\formal{second-order}
\scope The asserted remainder is $o(a^2)$. This statement supplies no
fourth-order remainder estimate.

\subsection{An arbitrary family approaching threshold}
\begin{claim}
Let $D>0$, let $\ell$ be a filter on an index type, and let $K_x,r_x$ be
real-valued functions on that type. Suppose $K_x\to2D$, $r_x\to0$ along
$\ell$, and eventually
\[
 K_x>2D,\qquad r_x>0,\qquad
 r_x=\operatorname{selfConsistency}(K_x,D,r_x).
\]
Then $r_x^2/(K_x-2D)\to1/D$ along $\ell$.
\end{claim}
\begin{proof}
Set $a_x=K_xr_x/D$ and write
\[
 E(a)=\frac12-\frac{a^2}{16}+a^2q(a).
\]
The previous expansion gives $q(a)\to0$ as $a\to0$, taking the quotient
with Lean's totalised division at zero. The assumed limits imply $a_x\to0$,
hence $q(a_x)\to0$.

For an index satisfying the eventual hypotheses, positivity gives
$a_x\ne0$ and $K_x-2D\ne0$. The previously proved reformulation of the
self-consistency equation gives $E(a_x)=D/K_x$. Substituting the expansion
and rearranging yields
\[
 K_x-2D=K_xa_x^2\bigl(1/8-2q(a_x)\bigr),\qquad
 \frac{r_x^2}{K_x-2D}
 =\frac{D^2}{K_x^3\bigl(1/8-2q(a_x)\bigr)}.
\]
The denominator in the last expression tends to
$(2D)^3/8=D^3>0$. The quotient therefore tends to $D^2/D^3=1/D$.
The equality of the quotients holds eventually, which transfers this limit
to the desired expression.
\end{proof}
\formal{family-exponent}
\scope Convergence to zero and eventual stationarity are hypotheses. This
general statement allows the bottom filter, where its eventual and convergence
conditions are vacuous. A concrete branch supplies the substantive application
below. The claim establishes stationary asymptotics, with no time evolution
or stability conclusion.

\subsection{A branch that satisfies the hypotheses}
For $D>0$ and $K>2D$, define $r_*(D,K)$ by choosing the coherent solution from
the theorem that proves existence and uniqueness in $(0,1]$.
The Lean function \texttt{coherentBranch} uses this choice in that regime
and returns zero outside it.

\begin{claim}
If $D>0$ and $K>2D$, then $0<r_*(D,K)\leq1$ and
\[
r_*(D,K)=\operatorname{selfConsistency}(K,D,r_*(D,K)).
\]
\end{claim}
\begin{proof}
The two inequalities on $D,K$ select the positive branch of the definition.
The specification of the chosen witness is exactly positivity, the bound by
one, and the fixed-point equation supplied by the existence theorem.
\end{proof}
\formal{branch-spec}

\begin{claim}
For fixed $D>0$, as $K$ decreases to $2D$ through $K>2D$,
\[
 \frac{r_*(D,K)^2}{K-2D}\longrightarrow\frac1D.
\]
\end{claim}
\begin{proof}
Apply the family theorem to the right-hand neighbourhood filter of $2D$,
with the identity function for $K$. The previously proved
\texttt{coherentBranch\_tendsto\_zero} supplies the limit of $r_*$ at the
threshold. Membership in the right-hand neighbourhood gives $K>2D$
eventually, and the branch specification supplies positivity and stationarity
there. These are all the hypotheses of the family theorem.
\end{proof}
\formal{branch-exponent}
\scope The approach filter on the real line is nontrivial, and the branch
provides actual coherent solutions. This removes the vacuity issue in this
application. The conclusion still concerns stationary solutions.

\section{What the conditional chain composes}
The chain uses a common topological and measured substrate, its sheaf of
global sections, a supplied complete metric on those sections, a neural field,
a thermodynamic cover $T$, and a reflexive boundary with prediction map $P$.
The precise type and instance parameters are printed in the appendix.

\subsection{The common downstream argument}
Suppose a sequence of energies coarse-grains to $L$, the elapsed parameter
$\tau$ is positive, and the following named implications hold:
\begin{itemize}
\item $E56$: coarse-graining implies $D>0$, $K\geq0$, $L=K$, and an
identification of the supplied field with the parameters $K,D$.
\item $E67$: those scalar conditions imply $K>2D$.
\item $E78$: a coherent order parameter implies that the supplied cover is
reached by the specified relaxation at a coupling at least $K$.
\item $E89$: that reachability supplies a section $s$ which restricts to the
local data of $T$, is fixed by $P$, and comes with the stated Lipschitz bound
on $P$ at \texttt{resonanceRate}$(K,D,\tau)$.
\end{itemize}

\begin{claim}
There is a global section compatible with a thermodynamic cover which is the
unique fixed point of the prediction map.
\end{claim}
\begin{proof}
Apply $E56$ to the coarse-graining hypothesis and $E67$ to its scalar
conclusion. The supercritical existence theorem supplies a coherent order
parameter. Apply $E78$ and then $E89$ to obtain the section $s$, its local
restrictions, the Lipschitz bound, and $P(s)=s$.

The section $s$ makes the space of global sections nonempty. The theorem
\nolinkurl{self_of_coherent_order_parameter}, using completeness, the
Lipschitz bound, coherence, and $\tau>0$, applies the contraction argument to
obtain a unique fixed point $p$. Since $s$ is fixed, uniqueness gives $s=p$.
Every fixed point $t$ equals $p$ and hence $s$. Pair $s$ with the supplied
cover $T$ and its known restriction equations. This is exactly the predicate
\texttt{UnifiedSelf}.
\end{proof}
\formal{downstream-chain}
\scope The identity between the glued section and a fixed section is supplied
in $E89$; the contraction argument proves uniqueness. The concrete-field
conjunct of $E56$ is retained as an input constraint but is not used by this
proof's downstream fixed-point argument. The source also establishes a local
\texttt{Unity} fact which the concluding construction does not reference.
Its constants can still occur in the unnormalised proof term; a syntactic
dependency list alone therefore does not identify minimal premises.

\subsection{The passive composition}
The passive chain additionally supplies a finite nonempty register, its
statistical-mechanics instance and update, a field and vacuum set, and measurable
carrier types for predictive information. Its remaining named implications are:
\begin{itemize}
\item $E12$: the capacity bound implies that the field leaves the vacuum set.
\item $E23$: leaving that set implies the register update is not surjective.
\item $E34$: dissipation supplies a predictive-dissipation structure whose
thermal scale and dissipated work match the register's temperature and heat,
and whose nonpredictive information is strictly positive.
\item $E45$: the predictive bound implies the chosen coarse-graining statement.
\end{itemize}

\begin{claim}
Under $E12$ through $E89$ and the stated structural hypotheses, the passive
chain concludes \texttt{UnifiedSelf} for the supplied reflexive boundary.
\end{claim}
\begin{proof}
Finiteness gives the capacity bound. Apply $E12$ to obtain vacuum departure,
then $E23$ to obtain a nonsurjective update. The proved dissipation theorem,
using the supplied statistical-mechanics structure, gives positive heat.
Apply $E34$ and retain its predictive-dissipation witness $R$. The bound
associated with $R$ gives the predictive node. Apply $E45$ to obtain
coarse-graining and then invoke the downstream claim with $E56$ through $E89$.
\end{proof}
\formal{passive-chain}
\scope The eight named implications remain hypotheses. Conditions inside
the statistical-mechanics, thermodynamic-cover, and other structures are
additional obligations. The equalities and positivity in $E34$ constrain
which witness is admitted, while the passage through $E45$ uses its
predictive structure. The conclusion has no predicate asserting experience
or biological implementation.

\section{A concrete nonconstant family and its trace}
On the three-site type with sites left, mid, and right, the development defines
the nonnegative mass profile
\[
 d_t=\bigl(\max(1+\cos t,0),\ 1,\ \max(2+\cos t,0)\bigr).
\]
The use of $\max$ expresses \texttt{Real.toNNReal}; these are nonnegative
spatial masses, without a normalisation requirement.

\begin{claim}
$d_0\ne d_\pi$.
\end{claim}
\begin{proof}
Suppose the profiles were equal. Evaluating at the left site would equate
$\max(1+\cos0,0)=2$ and $\max(1+\cos\pi,0)=0$, a contradiction.
\end{proof}
\formal{nonconstant-witness}
\scope This witnesses nonconstant dependence on phase for this particular
family. It does not assert that every thermodynamic cover has such dependence,
or that this family models a biological substrate.

\subsection{Before and after the actual tactics}\label{sec:witness-trace}
The following excerpt is generated from the saved elaboration data for this
proof. It selects the contradiction assumption, the evaluation at the left
site, and the simplification that closes the goal. The complete nested records,
including wrappers sharing source spans, remain in \nolinkurl{data/witness.json}.
\VerbatimInput{generated/witness-trace.txt}


\clearpage
\appendix
\section{Extracted statements and direct project dependencies}
\label{sec:formal}
The following text is generated from elaborated Lean expressions. Universe
parameter names, complete dependency lists including Mathlib, source excerpts,
and goal states are retained in the JSON snapshots. Pretty-printing uses Lean's
mathematical notation and can suppress implicit arguments at applications;
the quantified hypotheses of the declaration remain in the displayed type.
Dependency lists here show direct project constants in the proof term. They
are not a minimal set of assumptions or a transitive explanation of every lemma.
% Generated by proof_companion/render.sh from saved JSON; do not edit.
\subsection{sheaf-property}\label{decl:sheaf-property}
\noindent\nolinkurl{PhysicsOfConsciousness.probability_is_sheaf}\par
\noindent\href{../PhysicsOfConsciousness/Phase5_GlobalSection.lean}{Source, line 19}; 0 nested tactic records.\par
\VerbatimInput{generated/sheaf-property.statement.txt}
\noindent Direct project constants in the proof term:
\VerbatimInput{generated/sheaf-property.dependencies.txt}
\subsection{sheaf-gluing}\label{decl:sheaf-gluing}
\noindent\nolinkurl{PhysicsOfConsciousness.sheaf_glue_unique}\par
\noindent\href{../PhysicsOfConsciousness/Phase5_GlobalSection.lean}{Source, line 52}; 494 nested tactic records.\par
\VerbatimInput{generated/sheaf-gluing.statement.txt}
\noindent Direct project constants in the proof term:
\VerbatimInput{generated/sheaf-gluing.dependencies.txt}
\subsection{probability-gluing}\label{decl:probability-gluing}
\noindent\nolinkurl{PhysicsOfConsciousness.probability_glue_unique}\par
\noindent\href{../PhysicsOfConsciousness/Phase5_GlobalSection.lean}{Source, line 108}; 0 nested tactic records.\par
\VerbatimInput{generated/probability-gluing.statement.txt}
\noindent Direct project constants in the proof term:
\VerbatimInput{generated/probability-gluing.dependencies.txt}
\subsection{thermodynamic-gluing}\label{decl:thermodynamic-gluing}
\noindent\nolinkurl{PhysicsOfConsciousness.global_section_from_thermodynamics}\par
\noindent\href{../PhysicsOfConsciousness/Phase5_GlobalSection.lean}{Source, line 198}; 0 nested tactic records.\par
\VerbatimInput{generated/thermodynamic-gluing.statement.txt}
\noindent Direct project constants in the proof term:
\VerbatimInput{generated/thermodynamic-gluing.dependencies.txt}
\subsection{approximate-gluing}\label{decl:approximate-gluing}
\noindent\nolinkurl{PhysicsOfConsciousness.ApproximateGluing.select_close}\par
\noindent\href{../PhysicsOfConsciousness/Phase5_GlobalSection.lean}{Source, line 372}; 47 nested tactic records.\par
\VerbatimInput{generated/approximate-gluing.statement.txt}
\noindent Direct project constants in the proof term:
\VerbatimInput{generated/approximate-gluing.dependencies.txt}
\subsection{second-order}\label{decl:second-order}
\noindent\nolinkurl{PhysicsOfConsciousness.vonMisesSRatio_second_order}\par
\noindent\href{../PhysicsOfConsciousness/Phase8_CriticalExponent.lean}{Source, line 196}; 200 nested tactic records.\par
\VerbatimInput{generated/second-order.statement.txt}
\noindent Direct project constants in the proof term:
\VerbatimInput{generated/second-order.dependencies.txt}
\subsection{family-exponent}\label{decl:family-exponent}
\noindent\nolinkurl{PhysicsOfConsciousness.coherent_solution_critical_exponent}\par
\noindent\href{../PhysicsOfConsciousness/Phase8_CriticalExponent.lean}{Source, line 224}; 340 nested tactic records.\par
\VerbatimInput{generated/family-exponent.statement.txt}
\noindent Direct project constants in the proof term:
\VerbatimInput{generated/family-exponent.dependencies.txt}
\subsection{branch-spec}\label{decl:branch-spec}
\noindent\nolinkurl{PhysicsOfConsciousness.coherentBranch_spec}\par
\noindent\href{../PhysicsOfConsciousness/Phase8_CriticalExponent.lean}{Source, line 288}; 84 nested tactic records.\par
\VerbatimInput{generated/branch-spec.statement.txt}
\noindent Direct project constants in the proof term:
\VerbatimInput{generated/branch-spec.dependencies.txt}
\subsection{branch-exponent}\label{decl:branch-exponent}
\noindent\nolinkurl{PhysicsOfConsciousness.coherentBranch_critical_exponent}\par
\noindent\href{../PhysicsOfConsciousness/Phase8_CriticalExponent.lean}{Source, line 317}; 16 nested tactic records.\par
\VerbatimInput{generated/branch-exponent.statement.txt}
\noindent Direct project constants in the proof term:
\VerbatimInput{generated/branch-exponent.dependencies.txt}
\subsection{downstream-chain}\label{decl:downstream-chain}
\noindent\nolinkurl{PhysicsOfConsciousness.Chain.chain_from_coarseGrains}\par
\noindent\href{../PhysicsOfConsciousness/Chain.lean}{Source, line 796}; 91 nested tactic records.\par
\VerbatimInput{generated/downstream-chain.statement.txt}
\noindent Direct project constants in the proof term:
\VerbatimInput{generated/downstream-chain.dependencies.txt}
\subsection{passive-chain}\label{decl:passive-chain}
\noindent\nolinkurl{PhysicsOfConsciousness.Chain.chain}\par
\noindent\href{../PhysicsOfConsciousness/Chain.lean}{Source, line 858}; 66 nested tactic records.\par
\VerbatimInput{generated/passive-chain.statement.txt}
\noindent Direct project constants in the proof term:
\VerbatimInput{generated/passive-chain.dependencies.txt}
\subsection{nonconstant-witness}\label{decl:nonconstant-witness}
\noindent\nolinkurl{PhysicsOfConsciousness.Examples.richDensity_not_const}\par
\noindent\href{../PhysicsOfConsciousness/Examples/Phase5.lean}{Source, line 100}; 19 nested tactic records.\par
\VerbatimInput{generated/nonconstant-witness.statement.txt}
\noindent Direct project constants in the proof term:
\VerbatimInput{generated/nonconstant-witness.dependencies.txt}


\section{Two extracted physical obligations}
These projection types are extracted from the structures used by the gluing
results. They show the equilibrium and overlap-agreement obligations described
in Section~\ref{sec:gluing}. The snapshots contain the other encountered
structure fields as well. The exporter follows project statement types and
parent fields; it does not claim a complete semantic expansion through every
definition body or every Mathlib structure.
\VerbatimInput{generated/physical-fields.txt}
\end{document}
